\documentclass[11pt]{article}
\usepackage{amsmath,amssymb,amsthm}
\usepackage[margin=1in]{geometry}

\newtheorem{theorem}{Theorem}
\newtheorem{lemma}{Lemma}
\newcommand{\bits}{\{0,1\}}

\title{Laboratory V96: A Surplus-Component Avoider}
\author{NC0\_k-Avoid Laboratory}
\date{2026-08-10}

\begin{document}
\maketitle

Let $C:\bits^N\to\bits^{N+1}$ be an explicit local circuit.  Form the bipartite
incidence graph between input coordinates and output gates.  For a connected
component $K$, write $n_K$ and $m_K$ for its numbers of inputs and outputs.

\begin{lemma}[Positive-surplus component]
At least one component satisfies $m_K>n_K$.
\end{lemma}

\begin{proof}
The components partition the $N$ input vertices and the $N+1$ output vertices,
including isolated unused inputs and constant outputs.  Therefore
\[
 \sum_K(m_K-n_K)=(N+1)-N=1,
\]
so at least one summand is positive.
\end{proof}

Define
\[
 \rho(C)=\min\{n_K:m_K>n_K\}.
\]

\begin{theorem}[Comparison-free surplus-component avoidance]
There is a deterministic algorithm producing a word outside
$\operatorname{Range}(C)$ in time
\[
 O\!\left(2^{\rho(C)}\operatorname{poly}(N)\right).
\]
Consequently the algorithm is polynomial whenever $\rho(C)=O(\log N)$.
\end{theorem}

\begin{proof}
Choose a positive-surplus component $K$ with $n_K=\rho$.  Enumerate all
$2^\rho$ assignments to its input coordinates and evaluate the $m_K$ local
outputs.  The resulting local range $R_K$ satisfies
\[
 |R_K|\le2^\rho<2^{m_K}
\]
because $m_K>\rho$.

It is unnecessary to enumerate the full local output space.  Consider the
first $2^\rho+1$ distinct $m_K$-bit strings in lexicographic order.  Since
$R_K$ has at most $2^\rho$ members, one candidate $y_K$ is not in $R_K$.
Find such a candidate by lookup in the enumerated local range.

Construct a global output word whose restriction to the output vertices of
$K$ is $y_K$, filling all other output coordinates arbitrarily.  A preimage of
this global word would restrict to an input assignment of $K$ realizing
$y_K$, contradicting $y_K\notin R_K$.

The local enumeration and lookup cost $2^\rho\operatorname{poly}(N)$.
\end{proof}

The theorem uses neither exact child comparison nor the V92 canonical word.  It
is a constructive restricted control only: $\rho(C)$ can be linear for a large
connected support component, so no all-instance polynomial-time conclusion or
worst-case runtime improvement is claimed.

\end{document}
